\scanpage{23}
% archive foliation: 22
$S$ \inserted{\unsure{noeth.}} muni d'un Idéal \inserted{\unsure{qcoh.}} de carré nul $J$,
$Z\xrightarrow{\unsure{\varphi|v}}X$, $Y$ sur $S$
\[
S_0=V(J)
\]
\margin{d'où $\unsure{\varphi_0|v_0}:Z_0\unsure{\longrightarrow|\text{—}}X_0$, $Y_0$ sur $S_0$}
\[
u_0:X_0\longrightarrow Y_0\quad\text{tel que}\quad u_0\varphi_0=v_0.
\]
On va étudier l'obstruction à relever $u_0$ en $u:X\longrightarrow Y$
tel que $u\varphi=v$. \inserted{et l'obstruction analogue pour la \unsure{situation déduite}
à l'aide d'un hom $m:J\longrightarrow G$, ou après localisation\ldots}
À priori, \deleted{\unsure{donc l'obstruction à relever}}
\inserted{les relèvements de}
$u_0$ en un $u$ (sans condition supplémentaire $u\varphi=v$)
correspondent aux splittages d'une extension sur $P_0$
\[
0\longrightarrow\underline{O}_{\Gamma_0}\otimes_{S_0}J\longrightarrow E\longrightarrow\unsure{I|\mathcal{I}}_{\Gamma_0}\longrightarrow0.
\]
Les relèvements pour $v_0$ correspondent aux splittages
d'une extension \unsure{$E'$} sur $P'_0=Z_0\times_{S_0}Y_0$
\[
0\longrightarrow\underline{O}_{\Gamma'_0}\otimes_{S_0}J\longrightarrow E'\longrightarrow\unsure{I|\mathcal{I}}_{\Gamma'_0}\longrightarrow0.
\]
\margin{SGA 3 \uline{III} 4.2}
D'ailleurs, on constate qu'on a un homomorphisme
d'extensions, compatible avec $P'_0\longrightarrow P_0$
\diagram{
\begin{tikzcd}
0\arrow[r]&\underline{O}_{\Gamma_0}\otimes_{S_0}J\arrow[r]\arrow[d]&E\arrow[r]\arrow[d]&\unsure{I|\mathcal{I}}_{\Gamma_0}\arrow[r]\arrow[d]&0\\
0\arrow[r]&\underline{O}_{\Gamma'_0}\otimes_{S_0}J\arrow[r]&E'\arrow[r]&\unsure{I|\mathcal{I}}_{\Gamma'_0}\arrow[r]&0
\end{tikzcd}}
(en fait, $E'$ se déduit de $E$ par extension de
la base $\unsure{P'_4|P'_0}\longrightarrow\unsure{P_4|P_0}$ \unsure{en} \unsure{avec}
\diagram{\begin{tikzcd}&P_0\arrow[d]\\Z\arrow[r]&X\end{tikzcd}};
\unsure{\}} il
faut préciser un peu ]. \deleted{\unsure{Considérons, puis on}}
\inserted{D'ailleurs, la}
\deleted{\unsure{Ce qui dérive}} deuxième ligne est
munie d'un splittage, et les solutions
cherchées sont les splittages de la première
\uline{compatibles} avec celui donné pour la seconde.
